\subsection{Confronto dei termini nei limiti all'infinito}
\label{sec:esercizi-limiti-limiti-all-infinito}

\begin{exercise}[Rapporto di polinomi dello stesso grado]
\label{ex:limiti-foglio-2311-01}
\difficoltaesercizio{0}
Calcolare il limite
\[
  \lim_{x\to+\infty}\frac{3x^2-x+5}{2x^2+4x+1}.
\]
\rimandoteoria{sec:lez09-limite-infinito}
\end{exercise}
\begin{proof}[Soluzione dell'esercizio]
Dividendo numeratore e denominatore per $x^2$, si ottiene
\[
\begin{aligned}
  \lim_{x\to+\infty}\frac{3x^2-x+5}{2x^2+4x+1}
    &=\lim_{x\to+\infty}
      \frac{3-\frac1x+\frac5{x^2}}{2+\frac4x+\frac1{x^2}}\\
    &=\frac32.
\end{aligned}
\]
\end{proof}

\begin{exercise}[Denominatore di grado maggiore]
\label{ex:limiti-foglio-2311-02}
\difficoltaesercizio{0}
Calcolare il limite
\[
  \lim_{x\to+\infty}\frac{x^2+2x+5}{2x^3-3x^2+9}.
\]
\rimandoteoria{sec:lez09-limite-infinito}
\end{exercise}
\begin{proof}[Soluzione dell'esercizio]
Dividendo numeratore e denominatore per $x^3$, si ottiene
\[
\begin{aligned}
  \lim_{x\to+\infty}\frac{x^2+2x+5}{2x^3-3x^2+9}
    &=\lim_{x\to+\infty}
      \frac{\frac1x+\frac2{x^2}+\frac5{x^3}}{2-\frac3x+\frac9{x^3}}\\
    &=0.
\end{aligned}
\]
Il rapporto tende a $0$ da valori positivi.
\end{proof}

\begin{exercise}[Quoziente polinomiale a meno infinito]
\label{ex:limiti-foglio-2311-03}
\difficoltaesercizio{0}
Calcolare il limite
\[
  \lim_{x\to-\infty}\frac{x^2-2}{x}.
\]
\rimandoteoria{sec:lez08-limite-casi}
\end{exercise}
\begin{proof}[Soluzione dell'esercizio]
Semplificando il rapporto, si ottiene
\[
  \lim_{x\to-\infty}\frac{x^2-2}{x}
  =\lim_{x\to-\infty}\left(x-\frac2x\right)
  =-\infty.
\]
\end{proof}

\begin{exercise}[Logaritmo di un polinomio]
\label{ex:limiti-foglio-2311-12}
\difficoltaesercizio{0}
Calcolare il limite
\[
  \lim_{x\to+\infty}\log_{10}(x^2-3).
\]
\rimandoteoria{thm:lez15-limiti-esponenziale-logaritmo}
\end{exercise}
\begin{proof}[Soluzione dell'esercizio]
Poiché $x^2-3\to+\infty$ per $x\to+\infty$ e $10>1$,
per il limite del logaritmo segue
\[
  \lim_{x\to+\infty}\log_{10}(x^2-3)=+\infty.
\]
\end{proof}

\begin{exercise}[Radicale con termine infinitesimo]
\label{ex:limiti-foglio-2311-13}
\difficoltaesercizio{0}
Calcolare il limite
\[
  \lim_{x\to+\infty}\sqrt{1+\frac4x}.
\]
\rimandoteoria{thm:lez16-limite-sostituzione}
\end{exercise}
\begin{proof}[Soluzione dell'esercizio]
Poiché $1+4/x\to1$ per $x\to+\infty$,
per continuità della radice quadrata si ha
\[
  \lim_{x\to+\infty}\sqrt{1+\frac4x}=1.
\]
\end{proof}

\begin{exercise}[Radicale con quoziente a meno infinito]
\label{ex:limiti-foglio-2311-14}
\difficoltaesercizio{0}
Calcolare il limite
\[
  \lim_{x\to-\infty}\sqrt{\frac{1-x}{x^2}}.
\]
\rimandoteoria{thm:lez16-limite-sostituzione}
\end{exercise}
\begin{proof}[Soluzione dell'esercizio]
Per $x<0$ il radicando è positivo e
\[
  \lim_{x\to-\infty}\frac{1-x}{x^2}
  =\lim_{x\to-\infty}\left(\frac1{x^2}-\frac1x\right)=0.
\]
Per continuità della radice quadrata segue
\[
  \lim_{x\to-\infty}\sqrt{\frac{1-x}{x^2}}=0.
\]
\end{proof}

\begin{exercise}[Radice cubica di una quantità divergente]
\label{ex:limiti-foglio-2311-16}
\difficoltaesercizio{0}
Calcolare il limite
\[
  \lim_{x\to-\infty}\sqrt[3]{2-x}.
\]
\rimandoteoria{sec:lez08-limite-casi}
\end{exercise}
\begin{proof}[Soluzione dell'esercizio]
Poiché $2-x\to+\infty$ per $x\to-\infty$,
segue
\[
  \lim_{x\to-\infty}\sqrt[3]{2-x}=+\infty.
\]
\end{proof}

\begin{exercise}[Prodotto con un radicale convergente]
\label{ex:limiti-foglio-2311-17}
\difficoltaesercizio{0}
Calcolare il limite
\[
  \lim_{x\to-\infty}x\sqrt{3+\frac5x}.
\]
\rimandoteoria{thm:lez14-limite-prodotto}
\end{exercise}
\begin{proof}[Soluzione dell'esercizio]
Poiché
\[
  \lim_{x\to-\infty}\sqrt{3+\frac5x}=\sqrt3>0,
\]
e $x\to-\infty$, il prodotto diverge negativamente:
\[
  \lim_{x\to-\infty}x\sqrt{3+\frac5x}=-\infty.
\]
\end{proof}
